% The sample deck offered on the opening page. It says what Lectern is, how to use it and how it
% works, and uses what it describes: builds, speaker notes, an outline, links between slides and
% a TikZ picture that assembles. Build with `make` in this directory; the PDF is committed.
% The screenshots in img/ are taken by shots.mjs (see the Makefile).
\documentclass[aspectratio=169,14pt]{beamer}
\usetheme{default}
\usepackage{pdfpc}
\usepackage{tikz}
\usepackage{fancyvrb}
\usetikzlibrary{arrows.meta,positioning}
\definecolor{lectern}{HTML}{0736A8}
\definecolor{warm}{HTML}{C2410C}
\setbeamercolor{structure}{fg=lectern}
\setbeamercolor{frametitle}{fg=lectern}
\setbeamerfont{frametitle}{series=\bfseries}
\setbeamertemplate{navigation symbols}{}
\setbeamerfont{subtitle}{size=\normalsize,series=\mdseries}
\setbeamerfont{title}{size=\huge,series=\bfseries}
\setbeamertemplate{itemize items}[circle]
\setbeamertemplate{footline}{\hfill\usebeamercolor[fg]{structure}\footnotesize\insertframenumber\hspace{1em}\vspace{0.6em}}
\tikzset{
  box/.style={draw=lectern, very thick, rounded corners=3pt, minimum width=2.6cm, minimum height=1.1cm, align=center},
  flow/.style={-{Stealth[length=7pt]}, very thick, lectern},
}
\tikzset{
  dev/.style={box, minimum width=3.3cm, minimum height=1.2cm, font=\small, inner sep=3pt},
  both/.style={{Stealth[length=7pt]}-{Stealth[length=7pt]}, very thick, lectern},
}
\newcommand{\key}[1]{\textbf{#1}}
% The site's address, in one place: change it here when Lectern moves to a domain of its own.
\newcommand{\site}{\texttt{benlockwood.com/talks}}
% A screenshot with a thin rim, so its dark edge reads against the white slide.
\newcommand{\shot}[2]{\IfFileExists{img/#2.png}{{\setlength{\fboxsep}{0pt}\color{black!40}\fbox{\includegraphics[width=#1]{img/#2.png}}}}{\fbox{\parbox[c][3cm][c]{#1}{\centering #2}}}}
\newcommand{\goto}[2]{\hyperlink{#1}{\scalebox{1.4}{\beamergotobutton{#2}}}}
\newcommand{\return}[2]{\hyperlink{#1}{\scalebox{1.4}{\beamerreturnbutton{#2}}}}
\title{A tour of Lectern}
\subtitle{A full-featured, browser-based PDF presenter}
\author{A sample deck, made with Beamer and TikZ}
\date{}
\begin{document}

% The listings, kept outside their frames: a frame cannot hold a literal \end{frame}.
\begin{SaveVerbatim}{codenotes}
\usepackage{pdfpc}              % in the preamble

\begin{frame}{Results}
  Our main estimate.
  \pdfpcnote{Say *why* this matters.\\- then the caveat}
  \only<2>{\pdfpcnote{A note for build 2 only.}}
\end{frame}
\end{SaveVerbatim}
\begin{SaveVerbatim}{codelinks}
\section{Results}               % an entry in the outline

\hyperlink{proof}{\beamergotobutton{The proof}}

\appendix
\begin{frame}[label=proof]{Backup: the proof}
\end{SaveVerbatim}
\begin{SaveVerbatim}{codebeamer}
  \usepackage{pgfpages}
  \setbeameroption{show notes on second screen=right}
  ... \note{Say why this matters.}
\end{SaveVerbatim}

\begin{frame}
  \titlepage
  \pdfpcnote{Welcome. These are speaker notes: only the presenter view shows them, never the audience.\\\\Press the right arrow, or click the slide, to go on.}
\end{frame}

\begin{frame}{The core idea}
  \begin{itemize}
    \item \textbf{A full-featured PDF presenter.} Fast, with working in-PDF links and keyboard shortcuts for quick navigation.
    \item<2-> \textbf{A presenter view for the speaker,} as in PowerPoint or Keynote: notes, timer, outline, annotation tools.
      \begin{itemize}
        \item on a second screen of the same computer, or
        \item on a different device (laptop, tablet, phone), which is then a remote control too
      \end{itemize}
    \item<3-> \textbf{Entirely in the browser.} Nothing to install; it works on any computer with web access.
  \end{itemize}
  \only<1>{\pdfpcnote{Slides are drawn ahead of time, so the next one appears at once.\\\\Any PDF works: Beamer, or a PowerPoint or Keynote export.}}
  \only<2>{\pdfpcnote{You are probably reading this in the presenter view. The audience view shows the slide and nothing else.}}
  \only<3>{\pdfpcnote{That matters most on a seminar room's computer, where you can't install anything: it only needs a browser.}}
\end{frame}

\begin{frame}{Roadmap}
  \tableofcontents
  \pdfpcnote{These are the deck's sections, as written in the LaTeX source. The Outline tab beside Notes lists the same four, and a click there jumps to one.}
\end{frame}

\section{How to use it}

\begin{frame}{How to use it}
  \begin{enumerate}
    \item Go to \site{} and \textbf{open your PDF.}\\
      Your talk is now live, with a one-word \textbf{talk code}.
    \item<2-> On any other device, go to \site{} and \textbf{enter the code.}
    \item<3-> Each device chooses what it shows:
      \begin{itemize}
        \item \textbf{Audience view:} the slides, full screen
        \item \textbf{Presenter view:} notes, controls, timer
      \end{itemize}
    \item<4-> \textbf{Every device stays in step.} Advance on one, and all the others advance.
  \end{enumerate}
  \only<1>{\pdfpcnote{The code is shown on the page where the view is chosen, and under the connection button in the presenter view. It is never shown on the slides.}}
  \only<2-3>{\pdfpcnote{Any number of devices can join, and each can change its view at any time: Esc twice returns to the page where the view is chosen.}}
  \only<4>{\pdfpcnote{The laser pointer, drawings, black and white screens and the clock are shared in the same way.\\\\If the network drops, the computer showing the slides keeps working by its own keyboard or clicker.}}
\end{frame}

\begin{frame}{Audience view: the slides, full screen}
  \begin{columns}[c]
    \begin{column}{0.62\textwidth}
      \shot{\textwidth}{audience}
    \end{column}
    \begin{column}{0.36\textwidth}
      \begin{itemize}
        \item Shows the slides and nothing else.
        \item Runs a talk by itself: every key works here too.
        \item Move the mouse and click the Lectern mark, or tap twice, for the tools.
      \end{itemize}
    \end{column}
  \end{columns}
  \pdfpcnote{The picture shows the bar that the Lectern mark at the bottom left brings up: overview, laser, draw, black and white screens, the clock and Exit. It hides itself again.\\\\With one computer and one screen, this view is all you need.}
\end{frame}

\begin{frame}{Presenter view: notes, next slide, clock}
  \begin{columns}[c]
    \begin{column}{0.62\textwidth}
      \shot{\textwidth}{presenter}
    \end{column}
    \begin{column}{0.36\textwidth}
      \begin{itemize}
        \item On a second screen, a laptop, a tablet or a phone.
        \item Click the clock to set a countdown.
        \item \textbf{Outline} lists the sections; \textbf{Overview} shows every frame.
      \end{itemize}
    \end{column}
  \end{columns}
  \pdfpcnote{Laser and Draw work from here: point or write on the current slide, and the audience sees it.\\\\On a phone this view is mostly two large buttons, Prev and Next, with the notes above them.}
\end{frame}

\begin{frame}{Case 1: an iPad, and the room's computer}
  \begin{columns}[c]
    \begin{column}{0.4\textwidth}
      \centering
      \begin{tikzpicture}[node distance=0.75cm and 0.5cm]
        \node[dev] (pc) {Room computer\\\textbf{Audience view}};
        \node[dev, below=of pc] (ipad) {Your iPad\\\textbf{Presenter view}};
        \uncover<2->{\node[dev, below=of ipad, draw=warm] (phone) {Your phone,\\as a remote};}
        \draw[both] (ipad) -- node[right, font=\scriptsize, lectern] {in step} (pc);
        \uncover<2->{\draw[both, warm] (phone) -- (ipad);}
      \end{tikzpicture}
    \end{column}
    \begin{column}{0.58\textwidth}
      \begin{itemize}
        \item \textbf{Beforehand:} open your talk on the iPad and choose Presenter view.
        \item \textbf{In the room:} stand the iPad facing you. On the room's computer open \site, enter the code, choose Audience view.
        \item<2-> \textbf{Advance} with a clicker plugged into either one, or join on your phone and use that.
      \end{itemize}
    \end{column}
  \end{columns}
  \only<1>{\pdfpcnote{Nothing is installed or left behind on the room's computer. A private window there forgets the code as well.}}
  \only<2>{\pdfpcnote{A clicker sends arrow keys, so it works on whichever device it is plugged into.\\\\The phone joins with the same code and shows the presenter view: Prev, Next and your notes.}}
\end{frame}

\begin{frame}{Case 2: two displays on the room's computer}
  \begin{columns}[c]
    \begin{column}{0.4\textwidth}
      \centering
      \begin{tikzpicture}[node distance=0.75cm and 0.5cm]
        \node[dev] (proj) {Projector\\\textbf{Audience view}};
        \node[dev, below=of proj] (pc) {Room computer};
        \node[dev, below=of pc] (mon) {Its own screen\\\textbf{Presenter view}};
        \draw[flow] (pc) -- (proj);
        \draw[flow] (pc) -- (mon);
        \uncover<2->{\draw[both, warm] (proj.east) to[bend left=50] node[right, font=\scriptsize, warm] {swap} (mon.east);}
      \end{tikzpicture}
    \end{column}
    \begin{column}{0.58\textwidth}
      \begin{itemize}
        \item Bring the PDF on a USB drive, or download it there.
        \item Open \site{} and load the slides.
        \item Click \textbf{Both views:} the slides go to one display and the presenter view to the other.
        \item<2-> Wrong way round? Click \textbf{Swap displays.}
      \end{itemize}
    \end{column}
  \end{columns}
  \only<1>{\pdfpcnote{The displays must extend the desktop. Mirrored displays look like one screen to the browser.\\\\Chrome and Edge ask once for permission to place windows on the screens.}}
  \only<2>{\pdfpcnote{Swap displays is at the top of the presenter view.}}
\end{frame}

\begin{frame}{Case 3: no computer, phone to projector}
  \begin{columns}[c]
    \begin{column}{0.4\textwidth}
      \centering
      \begin{tikzpicture}[node distance=0.75cm and 0.5cm]
        \node[dev] (proj) {Projector};
        \node[dev, below=of proj] (phone) {Your phone\\\textbf{Audience view}};
        \node[dev, below=of phone] (ipad) {Your iPad\\\textbf{Presenter view}};
        \draw[flow, warm] (phone) -- node[right, font=\scriptsize, warm] {USB-C or HDMI} (proj);
        \draw[both] (ipad) -- node[right, font=\scriptsize, lectern] {in step} (phone);
      \end{tikzpicture}
    \end{column}
    \begin{column}{0.58\textwidth}
      \begin{itemize}
        \item Open your talk on the iPad and choose Presenter view.
        \item Join on your phone with the code and choose Audience view.
        \item Connect the phone to the projector with a USB-C or HDMI adapter.
      \end{itemize}
    \end{column}
  \end{columns}
  \pdfpcnote{The phone needs no computer in the room at all.\\\\Turn the phone sideways. On an iPhone, add the site to the Home Screen first: opened from there, the slides fill the screen without Safari's bars.}
\end{frame}

\begin{frame}{Case 4: no network, carry a packed copy}
  \begin{columns}[c]
    \begin{column}{0.56\textwidth}
      \shot{\textwidth}{pack}
    \end{column}
    \begin{column}{0.42\textwidth}
      \begin{itemize}
        \item \textbf{Download a packed copy} saves one HTML file: Lectern with your slides inside.
        \item Put it on a thumb drive and open it in Chrome or Edge.
        \item It needs no network, and shows the Audience view only, without a clock.
      \end{itemize}
    \end{column}
  \end{columns}
  \pdfpcnote{The button is on the page where the view is chosen, after a PDF is opened from a file.\\\\A packed copy shows the Audience view on its one computer: no presenter view, and no other devices can join it. Every key and the Menu bar still work there.}
\end{frame}

\section{How it works}

\begin{frame}{How it works}
  \centering
  \begin{tikzpicture}[node distance=1.2cm and 1.5cm, font=\small]
    \node[box] (mine) {Your device:\\the PDF};
    \uncover<2->{\node[box, right=of mine, draw=warm, minimum width=3.6cm] (relay) {The relay:\\a minimal chat room};}
    \uncover<2->{\node[box, right=of relay] (other) {A device that\\joins with the code};}
    \uncover<2->{\draw[both] (mine) -- (relay);}
    \uncover<2->{\draw[both] (relay) -- (other);}
    \uncover<3->{\node[below=0.15cm of relay, font=\footnotesize, warm] {sees only scrambled messages};}
  \end{tikzpicture}

  \medskip
  \begin{itemize}\small
    \item \textbf{Your slides are never stored on a server.} The PDF opens in your browser.
    \item<2-> Devices with the same code meet in a \textbf{chat room} that only passes their messages along: the slides, then where the talk is. It keeps nothing.
    \item<3-> Every message is \textbf{end-to-end encrypted} with a key made from the talk code, which the relay never sees.
  \end{itemize}
  \only<1>{\pdfpcnote{A TikZ picture that gains a piece with each build. The next-slide preview shows what the next key press adds.}}
  \only<2>{\pdfpcnote{Nothing leaves your device until another device joins with the code. Then the PDF is sent to it once, through the relay, and after that only small messages: frame 12, laser here, clock paused.}}
  \only<3>{\pdfpcnote{The code is the key, so anyone who has the code can join the talk. When the talk ends, every device forgets the slides.\\\\The notes travel too, and never reach an audience view.}}
\end{frame}

\section{Designing slides}

\begin{frame}{Any PDF works, but it is built for Beamer}
  \begin{itemize}
    \item \textbf{Any PDF opens,} including PowerPoint and Keynote exports.
    \item<2-> \textbf{Beamer decks get more:}
      \begin{itemize}
        \item overlays (\texttt{\textbackslash pause}, \texttt{<2->}) become builds of one frame
        \item \texttt{\textbackslash section} becomes the outline
        \item \texttt{\textbackslash hyperlink} buttons jump, and Back returns
        \item \texttt{\textbackslash pdfpcnote} becomes speaker notes
      \end{itemize}
  \end{itemize}
  \pdfpcnote{None of this needs a special class or theme: it is read from an ordinary Beamer PDF.\\\\The following frames show each of the four in turn.}
\end{frame}

\begin{frame}{Overlays become the builds of one frame}
  This frame is five PDF pages, and Lectern counts it as one frame.
  \begin{itemize}
    \item<2-> \key{→} and \key{←} step one build at a time.
    \item<3-> \key{Shift} with an arrow moves a whole frame.
    \item<4-> \key{A number, then Enter,} goes to that frame.
    \item<5-> \key{Tab} shows every frame at once.
  \end{itemize}
  \only<1>{\pdfpcnote{The frame counter stays at 13 while the five builds appear. Frame numbers match Beamer's own, bottom right.}}
  \only<2>{\pdfpcnote{Each build can carry a note of its own. This one belongs to the build about the arrow keys.}}
  \only<3>{\pdfpcnote{Try it: Shift and the right arrow skips the remaining builds.}}
  \only<4>{\pdfpcnote{Type 18 and press Enter to go to frame 18. Delete or Backspace returns here.}}
  \only<5>{\pdfpcnote{The overview keeps working while you present: the keys still move the talk.}}
\end{frame}

\begin{frame}[label=links]{Outline from \texttt{\textbackslash section}, links from \texttt{\textbackslash hyperlink}}
  \BUseVerbatim[fontsize=\small]{codelinks}

  \medskip
  Click a button to jump; \textbf{Back} (or \key{Delete}) returns.

  \smallskip
  \goto{backup-jump}{Try it: a backup slide}
  \pdfpcnote{The Outline tab beside Notes lists this deck's sections; a click on one jumps there.\\\\Links work with a click in either view.}
\end{frame}

\begin{frame}[label=notes]{Speaker notes: \texttt{\textbackslash pdfpcnote} inside the frame}
  \BUseVerbatim[fontsize=\small]{codenotes}

  \medskip
  The notes travel inside the PDF. They may use a little Markdown.

  \smallskip
  \goto{backup-notes}{Other ways to write notes}
  \pdfpcnote{This note is written in Markdown, as the notes of your own deck can be:\\\\- *emphasis* with asterisks\\- `code` with backticks\\- a list, one dash per line\\\\Two backslashes in the LaTeX source start a new line. A- and A+ above change the size of this text.}
\end{frame}

\section{Useful features}

\begin{frame}{Tools, on any device in the talk}
  \begin{itemize}
    \item \textbf{Laser:} point with the mouse or a finger; the audience sees the dot.
    \item<2-> \textbf{Draw:} write on the current slide, in three colours. The ink goes when the slide does.
    \item<3-> \textbf{Whiteboard:} \key{W} gives a white screen with Draw already on. \key{W} again returns to the slide.
    \item<4-> \textbf{Black:} \key{B} darkens the screen while you talk.
    \item<5-> \textbf{Clock:} counts up, or down from a length you set.
    \item<6-> \textbf{Overview:} \key{Tab} shows every frame; click one to go there.
  \end{itemize}
  \only<1-2>{\pdfpcnote{Each tool is a button in the presenter view and a key in either view.\\\\On the slides themselves, move the mouse and click the Lectern mark at the bottom left for the same tools.}}
  \only<3-4>{\pdfpcnote{What is written on the whiteboard is kept while you return to the slides, until it is cleared.}}
  \only<5-6>{\pdfpcnote{Click the clock to set a countdown. T pauses it.}}
\end{frame}

\begin{frame}{The room's computer keeps going if the network fails}
  \begin{itemize}
    \item When a device joins, the \textbf{whole PDF} is sent to it and held in that browser's memory.
    \item<2-> The page says so, under the preview: \textbf{Slides loaded locally.}
    \item<3-> From then on the slides are drawn on that computer. If the network fails, its keyboard or a clicker plugged into it still runs the talk.
    \item<4-> Other devices fall back into step by themselves when the connection returns.
  \end{itemize}
  \only<1-2>{\pdfpcnote{So load the talk on the room's computer before you start, and look for that line.\\\\Nothing is written to its disk: closing the tab, or End talk, forgets the slides.}}
  \only<3-4>{\pdfpcnote{Without a network a phone or tablet can't steer the room's computer, so know where its arrow keys are.\\\\For a room with no network at all, see Case 4, the packed copy.}}
\end{frame}

\begin{frame}{Keyboard shortcuts}
  \begin{columns}[T]
    \begin{column}{0.48\textwidth}
      \textbf{Moving}
      \begin{itemize}
        \item \key{→ ←} one build
        \item \key{Shift → ←} one frame
        \item \key{number, Enter} that frame
        \item \key{Delete} back after a jump
        \item \key{Tab} overview
      \end{itemize}
    \end{column}
    \begin{column}{0.48\textwidth}
      \textbf{Tools}
      \begin{itemize}
        \item \key{L} laser pointer
        \item \key{D} draw, \key{C} clear the ink
        \item \key{B} black, \key{W} whiteboard
        \item \key{T} pause the clock
        \item \key{?} every shortcut
      \end{itemize}
    \end{column}
  \end{columns}
  \pdfpcnote{A clicker works as it is: it sends the same keys.\\\\Esc twice leaves a view and returns to the page where the view is chosen.}
\end{frame}

\begin{frame}{iPhone or iPad: add to the Home Screen}
  \begin{enumerate}
    \item Open \site{} in \textbf{Safari.}
    \item Tap the \textbf{Share} button (the square with an arrow pointing up).
    \item<2-> Scroll down and choose \textbf{Add to Home Screen,} then tap \textbf{Add.}
    \item<3-> Open Lectern from its new icon, as you would an app.
  \end{enumerate}

  \medskip
  \uncover<3->{It then fills the whole screen, with no address bar or tabs.}
  \only<1-2>{\pdfpcnote{On an iPhone, Share may be behind the ... button at the bottom of Safari.\\\\Safari on an iPhone or iPad has no full screen for a web page, so this is the way to give the slides or the presenter view the whole screen there.}}
  \only<3>{\pdfpcnote{Worth doing once on each device you present from: the iPad for your notes, the phone as a remote or to feed a projector.\\\\It is still the web page: it needs the network to open a talk or join one.}}
\end{frame}

\begin{frame}{Now open your own deck}
  \begin{itemize}
    \item Press \key{Esc} twice to leave this view.
    \item \textbf{Leave this talk,} then open a PDF from this device.
    \item Its talk code lets your phone or laptop join.
  \end{itemize}
  \pdfpcnote{The last frame of the main talk. What follows is the appendix, reached by the buttons on frames 14 and 15.}
\end{frame}

\appendix
\section{Backup slides}

\begin{frame}[fragile,label=backup-notes]{Backup: other ways to write notes}
  \begin{itemize}
    \item \textbf{pdfcomment:} \verb|\pdfcomment{...}| inside a frame.
    \item \textbf{Beamer's own notes,} on a second half of each page:
  \end{itemize}
  \BUseVerbatim[fontsize=\small]{codebeamer}

  \medskip
  Lectern shows the slide half to the audience and the notes half to you.

  \smallskip
  \return{notes}{Back}
  \pdfpcnote{The i beside Notes in the presenter view lists all of these.}
\end{frame}

\begin{frame}[label=backup-jump]{Backup: you jumped here}
  This frame is in the appendix, after the last frame of the talk.
  \begin{itemize}
    \item<2-> The button below returns, as written in the deck.
    \item<2-> \textbf{Back}, under the current slide, returns from any jump.
  \end{itemize}

  \bigskip
  \return{links}{Back to the links}
  \pdfpcnote{A backup slide with two builds. The return button, the Back button and the Delete key all lead back.}
\end{frame}

\end{document}
